% === II. RELATED WORK ========================================================
\section{Related Work}
\label{sec:related}

\subsection{Tire modeling and Magic Formula identification}

The Magic Formula of Bakker, Nyborg and Pacejka~\cite{pacejka1987sae} and its
subsequent consolidation~\cite{pacejka1992magic,pacejka2012book} remains the
standard empirical description of steady-state tire force generation, and is
the model identified in this work. Its lateral form,
\begin{equation}
F_{y}=F_{z}\,D\sin\!\big(C\arctan\!\big(B\alpha-E(B\alpha-\arctan B\alpha)\big)\big),
\label{eq:mf}
\end{equation}
is parameterised by a stiffness factor $B$, shape factor $C$, peak factor $D$
and curvature factor $E$. Two properties of \eqref{eq:mf} are structural and
are used repeatedly in this paper. First, its slope at zero slip is
$F_{z}BCD$, so the axle cornering stiffness depends on the coefficients only
through that product and $E$ drops out. Second---and this is the identifiability
statement---if the data does not approach the peak, \eqref{eq:mf} linearises to
$F_{y}\approx F_{z}(BCD)\alpha$ and $B$, $C$, $D$ become mutually
unidentifiable. This is a standard persistent-excitation argument
\cite{ljung1999sysid}; we do not claim it as novel, but we do claim that the
baseline pipeline's virtual-data step can silently violate it, and that the
violation is what produced the failure we report.

Classical characterisation obtains the coefficients from steady-state
manoeuvres in a controlled space \cite{ontrack2024,amz2020}, which yields
clean data at high logistical cost. Direct fitting to measured contact-patch
forces, where a force-measuring hub or a simulator's own wheel telemetry is
available, is the most informative variant; in this work it is used only as a
\emph{reference} against which the on-track identifier is scored, never as
part of the identification loop.

\subsection{On-track and data-driven identification}

On-track methods trade excitation quality for operational simplicity. NLS
applied to lap data is the classical baseline and is known to be brittle
under noise \cite{ontrack2024}. Gaussian-process approaches learn the model
mismatch of a nominal model and have been used both for identification and
directly inside the controller
\cite{hewing2018cautious,hewing2020cautious,kabzan2019lbmpc}; they are
accurate but carry a continuing computational cost and, as
\cite{ontrack2024} observes, are usually evaluated starting from an already
accurate physical model. Neural approaches replace the GP with a cheaper
regressor at the cost of generalisation guarantees.

Deep Dynamics~\cite{deepdynamics2024} is the closest work on the
physics-constraint axis: it estimates Pacejka coefficients inside a neural
network and adds a \emph{Physics Guard} layer that clamps the internal
coefficient estimates into nominal physical ranges. Our admissibility gates
share that motivation but differ in placement and in what is constrained. We
found---and report as a substantive finding---that per-coefficient box
constraints are \emph{insufficient}: the degenerate fit that motivated this
work satisfied every box constraint while five of eight coefficients sat
exactly on a bound, and a later fit satisfied every box constraint while
being pairwise oversteering and open-loop unstable above
\SI{17.4}{\meter\per\second}. We therefore gate on derived quantities
(cornering stiffness, understeer sign, critical speed, grip ceiling versus
the trajectory's demand) at the model-to-controller interface, and treat a
coefficient sitting on a bound as a diagnostic of non-identifiability rather
than as a successful clamp.

The hybrid residual-plus-virtual-data structure of Dikici
\emph{et al.}~\cite{ontrack2024} is the baseline we extend, and is described
in detail in Section~\ref{sec:method}. Its key idea---use the learned model
only where it interpolates, and extract physically parameterised coefficients
from it---is retained unchanged. What we correct is the mechanism that
decides \emph{where} the learned model is queried.

\subsection{Overlap with concurrent work}

Wu \emph{et al.}~\cite{visionsysid2026} extend the same baseline pipeline
concurrently and independently, and the overlap must be stated plainly. They
introduce (i) a vision-based friction prior from a lightweight CNN over road
texture, used to warm-start the optimisation; (ii) an S4 structured
state-space model in place of the MLP, to capture temporal residuals; and
(iii) Nelder--Mead for the iterative Pacejka extraction, with virtual
simulation in CarSim. The pipeline described in this paper independently
contains an S4-family temporal residual architecture (an S4D
variant~\cite{gu2022s4d,gu2022s4} evaluated by recurrent scan rather than FFT
convolution, selected for the short $\approx$1500-sample sequences here), a
friction warm-start, and Nelder--Mead and differential-evolution solver
options alongside trust-region reflective.

We therefore do \emph{not} claim the temporal-residual architecture, the
friction warm-start concept, or the derivative-free solver as contributions
of this work. Two differences are worth recording because they are technical
rather than merely chronological. First, our friction prior is
proprioceptive---derived from measured accelerations, in the spirit of
model-free friction estimation from IMU data~\cite{slipfriction2025}---rather
than visual, and therefore does not require a camera or a textured-surface
assumption. Second, and more importantly, we find that a friction prior
computed as \emph{utilised} friction is a lower bound on \emph{available}
friction and cannot substitute for excitation: on a gentle lap the median
utilisation on our vehicle measures \num{0.044}. We consequently apply the
prior as a floor rather than as a replacement, and locate the actual fix in
the rollout excitation rather than in the initialisation. To our knowledge
the reachable-friction bound of Section~\ref{subsec:excitation} and its
consequence for the virtual-data step have not previously been stated.

\subsection{Learning-based and model predictive control for racing}

The downstream consumer of the identified model in this work is a nonlinear
MPC built on \texttt{acados}~\cite{acados2022}, using HPIPM partial
condensing~\cite{hpipm2020} and, in earlier configurations,
qpOASES~\cite{qpoases2014}. The path-tracking formulation follows standard
practice in racing MPC \cite{liniger2015,kabzan2019lbmpc}; the reference
raceline is generated by minimum-curvature optimisation in the manner of
Heilmeier \emph{et al.}~\cite{heilmeier2020}. Our contribution on this side
is not a new controller but the observation, made concrete in
Section~\ref{subsec:gates}, that an identified tire set can be numerically
inadmissible for the controller in ways the identifier itself has no way to
see---stiff linearisation, oversteering axle pairs, and grip ceilings below
the trajectory's demand---and that these are best enforced as an explicit
contract at the interface.

\subsection{Simulation-based validation and CARLA}

CARLA~\cite{carla2017} is the standard open simulator for autonomous-driving
research and provides the sensor and scenario infrastructure used here.
Its native vehicle dynamics come from the underlying PhysX wheeled-vehicle
model, which is a rigid-body model with per-wheel tire friction rather than
an empirical Magic Formula fit to measured data; this is the reason our study
treats the simulator as a well-instrumented \emph{plant} to be identified,
not as a source of ground-truth Pacejka coefficients. Recent work has
addressed exactly this fidelity gap: R-CARLA~\cite{rcarla2025} makes the
dynamics model interchangeable while retaining CARLA's sensor simulation and
reports substantial sim-to-real gap reductions, and Sagmeister
\emph{et al.}~\cite{simfidelity2026} quantify how far a vehicle model can be
simplified before conclusions about a trajectory-following controller stop
transferring---finding, relevantly for us, that the sensitivity is largest
near the acceleration limits. CARLA has also been used as the validation
environment for complete Formula Student Driverless stacks with ROS~2
interfaces \cite{carlaros2025}, and for digital-twin racing pipelines. The
platform described in Section~\ref{subsec:platform} sits in this line: it is
a ROS~2 bridge around a Formula Student vehicle model in CARLA, distinguished
by exposing the simulator's own per-wheel slip, load and force telemetry so
that an identification method can be scored against the forces the plant
actually produced.

Scaled autonomous racing platforms---F1TENTH~\cite{f1tenth2020} and the
ForzaETH Race Stack~\cite{forzaeth2024}---provide the software context of the
baseline. Full-scale Formula Student Driverless systems~\cite{amz2020}
provide the vehicle class targeted here.

\subsection{Gap addressed}

Summarising: the baseline~\cite{ontrack2024} establishes that a residual
network plus virtual steady-state data can recover a Pacejka model from a
lap, at 1:10 scale. Concurrent work~\cite{visionsysid2026} improves its
residual architecture and its initialisation. Neither addresses the
condition under which the virtual data itself carries information about the
tire's peak, which is the binding constraint at full scale. Neither treats
the admissibility of the identified model for the controller that will
consume it as an explicit, physically derived contract. This paper addresses
both, and does so on a platform where the answer can be checked against the
plant's own forces.
